\section{Introduction}
\label{sec:introduction}

A nef divisor has nonnegative degree on every curve. A semiample
divisor has a positive Cartier multiple generated by global sections,
and therefore defines a morphism to a projective variety. The passage
from numerical positivity to such a morphism is central to the minimal
model program. For an ordinary log canonical divisor, this is the
abundance problem. Adding an arbitrary nef divisor changes the natural
conclusion: one must allow replacement by a numerically equivalent
divisor.

We call a rational Cartier divisor \emph{numerically semiample} if it is
numerically equivalent to a semiample rational Cartier divisor on the
same variety. Here numerical equivalence means equality of degrees on
all integral curves. We prove the following statement.

\begin{theorem}\label{thm:main}
Let $(X,B)$ be a projective klt $\Q$-pair over an algebraically closed
field of characteristic zero. Suppose that $K_X+B$ is pseudo-effective,
and let $M$ be a nef $\Q$-Cartier divisor on $X$. If
\[
                         D=K_X+B+M
\]
is nef, then there is a semiample $\Q$-Cartier divisor $L$ on $X$ such
that $D\equiv L$.
\end{theorem}

The rationality assumptions apply to the boundary and the nef summand;
in particular the latter is a divisor on $X$ itself. The proof uses the
log-abundance theorem of \cite[Theorem~1.1]{LA} and the terminating
minimal-model programs of \cite[Theorem~1.1 and Section~3.3]{MM}.
Their precise forms are stated in Section~\ref{sec:preliminaries}.

Theorem~\ref{thm:main} contains the Generalised Abundance Conjecture
formulated by Lazi\'c--Peternell \cite{LP}, where the nef summand is
Cartier. Numerical equivalence is necessary even in dimension one.
Indeed, a nontorsion degree-zero line bundle on a complex elliptic curve
is nef and numerically trivial, but none of its positive powers has a
nonzero section. Since the canonical bundle of the curve is trivial,
this is an example with $B=0$ in Theorem~\ref{thm:main}.
For the ordinary adjoint, Fukuda's theorem
\cite[Theorem~0.1]{Fukuda} shows that a semiample numerical
representative makes $K_X+B$ itself semiample. For $M=0$, the conclusion
is therefore equivalent to ordinary klt abundance, which is an input
to this proof.

\begin{corollary}\label{cor:cy}
Let $(X,B)$ be a projective klt $\Q$-pair in characteristic zero with
$K_X+B\equiv0$. Every nef $\Q$-Cartier divisor on $X$ is numerically
semiample.
\end{corollary}
\begin{proof}
For a nef divisor $M$, the divisor $K_X+B+M$ is nef and numerically
equivalent to $M$. Apply Theorem~\ref{thm:main}.
\end{proof}

\subsection{Background and the role of the present argument}

The generalized-pair formalism retains a nef divisor on a higher
birational model as part of the adjoint data. Birkar--Zhang \cite{BZ}
developed effective birationality for these polarized pairs, extending
the tools used to study Iitaka fibrations. The canonical bundle formula
of Ambro \cite{Ambro} is one source of such nef data. In this paper the
original nef summand descends to $X$, but a minimal-model program forces
us to keep track of it as a fixed nef divisor on a higher model.

Lazi\'c--Peternell \cite{LP} formulated generalized abundance as a
numerical extension of abundance and related it to semiampleness on
varieties with trivial canonical class. Their Theorem~B, assuming
termination of flips and abundance in dimensions at most $n$, proves
numerical semiampleness when the ordinary adjoint has positive numerical
dimension. Its numerical-dimension-zero case also assumes their
Semiampleness Conjecture. Their Corollaries~C and~D give unconditional
results for surfaces and for threefolds with ordinary adjoint of positive
numerical dimension. These results identify the main remaining
positivity issue: an arbitrary nef summand on a variety with trivial
canonical class need not supply sections directly.

Subsequent work has addressed numerical nonvanishing, which asks for
an effective numerical representative, and refinements of numerical
semiampleness. Lazi\'c--Peternell \cite[Theorems~B and~C]{LPII}
obtained numerical nonvanishing under additional numerical or metric
hypotheses and, in the metric case, conjectural MMP and abundance
assumptions. For a Cartier nef summand, Chaudhuri
\cite[Corollary~3]{Chaudhuri} proved generalized abundance when the
nef-reduction base of $D$ has dimension two and $K_X+B$ is effective,
or when that base has dimension three and $\kappa(K_X+B)>0$.
Here nef reduction contracts the curves of degree zero through very
general points; its precise properties are recalled in
Section~\ref{sec:preliminaries}. On surfaces, Fontanari
\cite[Theorem~1.4]{Fontanari} proved actual semiampleness when the nef
summand is Cartier and its numerical class is not proportional to
that of the ordinary adjoint.

The birational reductions below follow \cite[Sections~4--7]{LP}.
The companion results \cite{LA,MM} supply ordinary log abundance and
the existence of the particular terminating programs required in those
reductions. The further ingredient proved here is that a nef divisor
of full nef dimension on a klt Calabi--Yau pair is big. Full nef
dimension means that every curve through a very general point has
positive degree. Bigness is stronger: it requires the volume of the
divisor to be positive. The gap between these two conditions is the
geometric issue addressed by the proof.

Our treatment adapts the moving-jet and Frobenius arguments of
\cite[Sections~7--9]{LA}. The moving-point method has antecedents in
the local-positivity work of Ein--K\"uchle--Lazarsfeld \cite{EKL}.
The finite-data comparison uses Frobenius jets, as studied by
Musta\c t\u a--Schwede \cite{MustataSchwede}, the canonical filtrations
of Sun and Kitadai--Sumihiro \cite{Sun,KS}, and Langer's control of
Frobenius slopes \cite{Langer}. We reproduce the comparison proof
with its uniformity conditions. The adaptation to an arbitrary nef
summand occurs in the geometric exclusions preceding it: the polarized
birationality theorem of Birkar--Zhang gives curves of bounded degree,
contradicting a growing integral lower bound for their degrees. We
give explicit proofs of those exclusions, including the case of
correspondences between finite covers.

\subsection{Structure of the proof}

There are two simultaneous induction statements. The first says that
$K_X+B+cM$ is big for every rational $c>0$ whenever one positive
combination of $K_X+B$ and $M$ has positive degree on all curves through
a very general point. The second says that the positive part of
$K_X+B+M$ on a suitable smooth birational model is rational and
numerically semiample. The positive part is the divisor left after
removing the asymptotically fixed prime divisors. When the original
adjoint is nef, it has no such negative part, and the second assertion
descends to Theorem~\ref{thm:main}.

Section~\ref{sec:reduction} sets up this induction. A pair of compatible
minimal-model programs retains both a semiample ordinary adjoint and a
nearby nef generalized adjoint. Positive Iitaka dimension then permits
induction on the fibres. The remaining full-dimension case is a klt
Calabi--Yau pair.

Sections~\ref{sec:jet-tools} and~\ref{sec:geometry} treat the essential
terminal case with trivial canonical bundle. If a nef Cartier divisor
$L$ of full nef dimension were not big, its ample approximations could
be rescaled to bounded volume while their degrees on very general
curves tend to infinity. The bounded-degree obstruction controls the
base loci of moving jet kernels. It produces a small vanishing-order
bound on the variety itself and a large jet system on a projective
bundle over its square. Section~\ref{sec:frobenius} shows that these two
conclusions are incompatible. The comparison is made on fixed complex
data before reduction to positive characteristic.

Section~\ref{sec:completion} passes through an index-one cover and a
boundary perturbation to obtain the full-dimension assertion for all
klt Calabi--Yau pairs. Section~\ref{sec:positive} proves the
positive-part assertion by nef reduction and the canonical bundle
formula, then descends numerical semiampleness to the original variety
and to any algebraically closed field of characteristic zero.
